\documentclass[12pt]{article}
\usepackage[a4paper,margin=2.5cm]{geometry}
\usepackage[T1]{fontenc}
\usepackage[utf8]{inputenc}
\usepackage{lmodern}
\usepackage{microtype}
\usepackage{needspace}
\usepackage{booktabs}
\usepackage{amsmath}
\usepackage{xcolor}
\usepackage{listings}
\usepackage[hidelinks]{hyperref}
\definecolor{codegray}{RGB}{245,245,245}
\lstset{basicstyle=\ttfamily\small,backgroundcolor=\color{codegray},frame=single,
  breaklines=true,showstringspaces=false,columns=fullflexible,keepspaces=true}
\setlength{\parindent}{0pt}
\setlength{\parskip}{0.55em}

\begin{document}
\pagenumbering{gobble}
\begin{titlepage}
\centering
{\large Ministerul Educa\c{t}iei \c{s}i Cercet\u{a}rii al Republicii Moldova\par}
\vspace{0.45cm}
{\large Universitatea Tehnic\u{a} a Moldovei\par}
\vspace{0.45cm}
{\large Facultatea Calculatoare, Informatic\u{a} \c{s}i Microelectronic\u{a}\par}
\vspace{0.45cm}
{\large Departamentul Ingineria Software \c{s}i Automatic\u{a}\par}
\vspace{3cm}
{\Large\bfseries Report on Laboratory Work No. 11\par}
\vspace{0.55cm}
{\LARGE\bfseries The Race and the Semaphore\par}
\vspace{0.35cm}
{\large Operating Systems\par}
\vfill
\begin{flushleft}
\textbf{Done by:} st. gr. FAF-241, Ca\c{t} Chiril\\[0.8cm]
\textbf{Checked by:} asist. univ. Rei\c{t}man Patricia
\end{flushleft}
\vfill
{Chi\c{s}in\u{a}u -- 2026\par}
\end{titlepage}
\pagenumbering{roman}
\tableofcontents
\clearpage
\pagenumbering{arabic}

\section{Aim, environment, and deliverables}
This lab investigates what happens when two POSIX threads update a shared
variable without synchronization. I first observed lost updates, then
protected the same increment with a mutex and a semaphore. The remaining
experiments use a semaphore to signal readiness and to limit the number of
cars admitted to a parking lot. All four required parts are complete; the
atomic-counter comparison is an additional, optional experiment.

The programs ran on Ubuntu 24.04.5 LTS, on aarch64 with glibc 2.39 and
GCC 13.3.0. Every build explicitly used \texttt{-O0}, as the handout asks
not to optimize the race example. The full flags were:
\begin{lstlisting}
gcc -std=c11 -Wall -Wextra -Werror -O0 -pthread \
    -o /tmp/race lab11/code/race.c
\end{lstlisting}
All nine configurations compiled successfully under these flags. The same
flags were used for the five required programs, the three altered-value
variants, and the optional atomic program. Each configuration ran three
times, giving \textbf{27 executions}. GNU time recorded elapsed, user,
and system time; the runner also measured end-to-end wall time with Python's
monotonic high-resolution clock.

The required sources are \texttt{race.c}, \texttt{mutex.c}, \texttt{sem.c},
\texttt{order.c}, and \texttt{parking.c}, all under \texttt{lab11/code/}.
They use the companion \texttt{thread\_utils.h} for consistent error
checks. POSIX thread functions return an error code rather than setting
\texttt{errno}; the helper handles that correctly. Semaphore waits retry
\texttt{EINTR}, so an interrupted wait never admits a thread without a
permit. Synchronization objects are destroyed after the worker threads
have joined. The race itself remains deliberately unprotected.

The submitted defaults are the correct versions: one permit for the
counter semaphore, zero for the readiness semaphore, and three parking
places. Separate builds with \texttt{-DSEM\_START=2},
\texttt{-DREADY\_START=1}, and \texttt{-DPLACES=1} implement the handout's
variations without leaving the final source in a broken configuration.
Parking IDs are passed as pointers to stable array elements rather than
integer-to-pointer casts.

The required \texttt{report.md} contains all ten Observe answers and
selected output. The PDF follows the previous laboratory reports; its
editable source and the handout are in \texttt{report\_files/}. Actual
outputs, timing records, and verification results are retained under
\path{lab11/report_files/output/}. No sample counter values from the
handout were substituted for measurements.

\clearpage
\section{Part 1: seeing the race}
\subsection{The shared increment}
Both threads execute one million increments of the same global
\texttt{long counter}, initially zero. Main joins both threads before
printing, so it does not read the final result while a worker is still
running. Joining does not, however, synchronize the two workers' updates
with each other.
\begin{lstlisting}[language=C]
for (int i = 0; i < N; i++)
    counter++;  /* Intentionally not protected. */
\end{lstlisting}
The expected mathematical total is \(2 \times 1000000 = 2000000\).
The observed totals were:
\lstinputlisting{output/race.txt}

\subsection{How one addition is lost}
An ordinary increment is conceptually a load, an addition in a register,
and a store back to memory. A possible overlap is:
\begin{center}
\begin{tabular}{lll}
\toprule
Step & Thread A & Thread B \\
\midrule
1 & Load counter = 7 & \\
2 & & Load counter = 7 \\
3 & Compute 7 + 1 = 8 & \\
4 & & Compute 7 + 1 = 8 \\
5 & Store 8 & \\
6 & & Store 8 \\
\bottomrule
\end{tabular}
\end{center}
Two attempted additions leave 8, not 9. The second store overwrites rather
than combines the first update. Scheduler decisions and overlapping
accesses explain why the measured totals differ between these runs.

This is a teaching model of a data race, not a valid way to compute a
counter. In the C memory model, conflicting unsynchronized accesses have
\textbf{undefined behavior}. Building with \texttt{-O0} makes the example
easier to observe but does not remove that undefined behavior. Contrary to
a literal reading of the handout's ``never 2,000,000'', a racy run can
accidentally print the expected number. Neither a fixed range of wrong
answers nor a different answer on every invocation is guaranteed.

\clearpage
\section{Part 2: mutex and binary semaphore}
The mutex version surrounds the complete increment with lock and unlock.
The semaphore version waits for one permit and returns it after incrementing:
\begin{lstlisting}[language=C]
/* mutex.c */
pthread_mutex_lock(&lock);
counter++;
pthread_mutex_unlock(&lock);
/* sem.c: s initially contains one permit */
sem_wait_checked(&s);
counter++;
sem_post(&s);
\end{lstlisting}
The excerpt omits error-check wrappers around the mutex and post operations.
Only one thread may update the counter at a time in either correct version.
The results of the three runs were:
\begin{center}
\begin{tabular}{lrrr}
\toprule
Configuration & Run 1 & Run 2 & Run 3 \\
\midrule
Mutex & 2000000 & 2000000 & 2000000 \\
Semaphore, start 1 & 2000000 & 2000000 & 2000000 \\
Semaphore, start 2 & 1791203 & 1709203 & 1844541 \\
\bottomrule
\end{tabular}
\end{center}
With two permits, both workers can enter simultaneously. The semaphore
protects its own count, not the application's unprotected increment; mutual
exclusion is lost and the altered program again has a data race. A binary
semaphore works as a lock here, but unlike a mutex it has no owner and can
be incorrectly posted above one.

\subsection{The price of synchronization}
Measured end-to-end wall time in seconds:
\begin{center}
\begin{tabular}{lrrrr}
\toprule
Program & Run 1 & Run 2 & Run 3 & Median \\
\midrule
race & 0.004462 & 0.004463 & 0.004460 & 0.004462 \\
mutex & 0.055737 & 0.056001 & 0.058797 & 0.056001 \\
sem & 0.131739 & 0.172545 & 0.142064 & 0.142064 \\
atomic & 0.049935 & 0.036728 & 0.031808 & 0.036728 \\
\bottomrule
\end{tabular}
\end{center}
These values include process startup and the GNU time wrapper. Separate
GNU elapsed/user/system values are retained in the raw timing files.
GNU time rounded the race's elapsed time to 0.00 s, not literally zero.

Two million lock/unlock or wait/post pairs add synchronization and shared
cache-line overhead; contention can also block and wake threads. The
measurements compare these implementations, including their error checks,
not isolated primitive latency. The observed median order was race,
atomic, mutex, semaphore, from fastest to slowest. Exact timings and the
relative order are not universal, and a fast incorrect answer is not a
valid optimization.

\clearpage
\section{Part 3: a semaphore as a signal}
\subsection{Starting at zero: wait until the data is ready}
The shared array \texttt{message[32]} is global. The writer simulates work
with a one-second sleep, copies \texttt{"hello"} into the array, announces
completion, and posts \texttt{ready}. The reader waits before reading:
\begin{lstlisting}[language=C]
/* Writer */
strcpy(message, "hello");
printf("A: message written\n");
sem_post(&ready);
/* Reader */
sem_wait_checked(&ready);
printf("B: I read \"%s\"\n", message);
\end{lstlisting}
With \texttt{sem\_init(\&ready, 0, 0)}, all three runs produced this output:
\lstinputlisting{output/order-1.stdout}
The successful wait consumes the permit posted after the write and
synchronizes the reader with the writer's prior memory operations. It is
the semaphore, not the sleep or the order of thread creation, that
establishes the required data dependency. Posting before the reader reaches
wait also works: the permit is retained until consumed.

\subsection{Starting at one: the broken ordering}
I compiled a separate version with \texttt{-DREADY\_START=1}. Every altered
run produced:
\lstinputlisting{output/order_one-1.stdout}
The reader immediately consumes the initial permit and proceeds before the
writer posts. In these runs it reads the empty, zero-initialized global
array while the writer is sleeping. This is not uninitialized garbage.
The altered program no longer guarantees write-before-read ordering;
other schedules may behave differently. A sleep is not a substitute for
synchronization of shared data.

Both configurations took about one second overall because main joins the
writer even when the reader finishes early. The submitted default remains
zero: no permission to proceed exists until a signal is posted.

\clearpage
\section{Part 4: the parking lot}
The semaphore starts at three and controls admission for six cars. Each car
waits before entering and posts only after leaving. A separate mutex keeps
the occupancy update and its print together. The mutex is released while
the car is parked, so it is not what limits the number of parked cars.
\lstinputlisting{output/parking-1.stdout}
The maximum reported occupancy was \textbf{3} in each of three default runs.
The guarantee comes from the wait before admission, the three initial
permits, and exactly one post per departure. No fixed arrival order or
FIFO fairness is promised. The verifier tracks the active car IDs and
checks every displayed occupancy, rather than merely searching for a
maximum number in the output.

\subsection{Changing the capacity to one}
The separately compiled \texttt{-DPLACES=1} version had a maximum occupancy
of \textbf{1} in all three runs. It behaves as a binary semaphore admitting
one car at a time for the entire parking interval. Every car still entered
and left exactly once, and the final occupancy was zero.

Default runs took approximately 2.01 seconds: roughly two waves of three
one-second stays. The one-place runs took approximately 6.02 seconds:
six sequential stays. These times are observations, not timing assertions.
The semaphore becomes a one-at-a-time admission lock, not a mutex with
ownership rules. The mutex remains useful for the occupancy diagnostics,
including the interval between releasing it and posting the permit.

\clearpage
\section{Handout questions and answers}
\subsection{Part 1: the race}
\begin{enumerate}
\Needspace{6\baselineskip}
\item \textbf{Question:} Counter++ looks like one step. Write the three steps the CPU really does.
\par\textbf{Answer:} Conceptually load the shared value into a register,
add one to the register, then store it back. This ordinary C operation is
not an atomic read-modify-write.

\Needspace{6\baselineskip}
\item \textbf{Question:} Describe, with those steps, how two threads can lose one addition.
\par\textbf{Answer:} Both load 7 before either stores; both compute 8;
each stores 8. The final value is 8 instead of 9, so one increment is lost.

\Needspace{6\baselineskip}
\item \textbf{Question:} Why is the result different at every run?
\par\textbf{Answer:} Scheduling and overlap of memory accesses change
which updates collide. All three recorded values differed. A C data race
has undefined behavior, so different results on every run are not actually
guaranteed; a correct-looking result is also possible.
\end{enumerate}

\subsection{Part 2: mutex and semaphore}
\begin{enumerate}
\Needspace{6\baselineskip}
\item \textbf{Question:} Run mutex and sem three times each. Is the answer always 2,000,000?
\par\textbf{Answer:} Yes, all six recorded runs returned exactly 2,000,000.
Both versions serialize the full increment, and main joins both threads.

\Needspace{7\baselineskip}
\item \textbf{Question:} What would happen in sem.c with sem\_init(\&s, 0, 2)? Try it and explain.
\par\textbf{Answer:} Two permits allow both workers into the critical
section. The observed totals were 1791203, 1709203, and 1844541. The
increment is no longer mutually exclusive and has a data race again.

\Needspace{6\baselineskip}
\item \textbf{Question:} Why is the locked version slower?
\par\textbf{Answer:} Every increment pays for lock/unlock or wait/post,
shared-cache-line contention, and possibly blocking and waking. The race
is faster here but does not compute a valid result.
\end{enumerate}

\subsection{Part 3: signaling}
\begin{enumerate}
\Needspace{7\baselineskip}
\item \textbf{Question:} Change the start value to 1, compile, run. What does B read now, and why?
\par\textbf{Answer:} In all three runs B read \texttt{""}, the initially
empty global buffer, before A wrote the message. The initial permit lets
B pass wait immediately, removing the required write-before-read ordering.

\Needspace{6\baselineskip}
\item \textbf{Question:} Put it back to 0. In one sentence: what does a semaphore that starts at 0 mean?
\par\textbf{Answer:} There is no permission to proceed yet, so a waiter
must wait for another thread's post, or consume a signal already posted.
\end{enumerate}

\subsection{Part 4: parking}
\begin{enumerate}
\Needspace{7\baselineskip}
\item \textbf{Question:} What is the largest value of inside you see? Which line of code guarantees it?
\par\textbf{Answer:} Three. \texttt{sem\_wait(\&places)}, called through
\texttt{sem\_wait\_checked}, before entry enforces the limit, with three
initial permits and one balanced post after each departure. The mutex
protects the diagnostic counter, not the full parking interval.

\Needspace{6\baselineskip}
\item \textbf{Question:} Set PLACES to 1. What does the semaphore become?
\par\textbf{Answer:} A binary semaphore acting as a one-at-a-time
admission lock. Each altered run had maximum occupancy one and final
occupancy zero; this does not give it mutex ownership or FIFO semantics.
\end{enumerate}

\clearpage
\section{Optional atomic counter and verification}
\subsection{An indivisible counter update}
The optional \texttt{atomic.c} uses \texttt{atomic\_long} and
\texttt{atomic\_fetch\_add(\&counter, 1)}. It returned exactly 2,000,000 in
all three runs, without an explicit mutex or semaphore. Its median elapsed
time was 0.036728 seconds in this measurement, below the mutex and semaphore
versions but above the incorrect race.

Atomic operations are indivisible at the C memory-model level. C does not
promise that every atomic type on every platform uses a single hardware
instruction or is lock-free; a target may use exclusive-access retries or
other implementation support. This solves one counter update, not an
arbitrary multi-variable critical section.

\subsection{What was verified}
The retained verifier passed \textbf{389 assertions across 27 executions}:
\begin{itemize}
\item All runs exited successfully without error diagnostics; raw standard
output matched the corresponding machine-readable record.
\item Mutex, semaphore-start-one, and atomic runs all reached exactly
2,000,000. Both intentionally unsafe counter configurations exhibited
incorrect totals in the recorded runs.
\item The zero-start reader printed \texttt{hello} only after the writer's
announcement. The altered-start experiment also completed and its actual
output was retained without pretending it was a correctness guarantee.
\item Every parking run contained six entries and six departures. Each car
entered and left exactly once. Every occupancy value matched the active
set, stayed within capacity, and returned to zero.
\end{itemize}
The tests impose no arbitrary timing threshold or fixed car order, and they
do not prove the absence of every possible concurrency bug. The safety
argument comes from the synchronization design; repeated runs provide
supporting evidence. All sources retained \texttt{-O0}; no artificial
counter result or forced sleep was inserted into the increment loops.

\section{Conclusion}
The unprotected increment loses work because its read-modify-write is not
atomic. A mutex or a one-permit semaphore serializes the update. A
zero-permit semaphore instead expresses readiness: the reader cannot
proceed until the writer publishes the data. A multi-permit semaphore
limits concurrent access to a finite resource, while a mutex can separately
protect diagnostic state. The experiments demonstrate three semaphore
roles---lock, signal, and resource counter---and show why correctness must
be established before comparing execution speed.
\end{document}
